\subsection{مبانی نظری و فرمول‌بندی ریاضی}
برای ماتریس وزن اولیه $\mathbf{W} \in \mathbb{R}^{d_{\text{out}} \times d_{\text{in}}}$، تجزیه مقادیر منفرد فاز نخست اجرا می‌شود:
\begin{equation}
\mathbf{W} = \mathbf{U} \mathbf{\Sigma} \mathbf{V}^\top
\end{equation}
با گزینش رتبه پایه $r \le \min(d_{\text{out}}, d_{\text{in}})$، تقریب کم‌رتبه $\mathbf{L}$ از ضرب فاکتورهای برش‌یافته شکل می‌گیرد:
\begin{equation}
\mathbf{L} = \mathbf{U}_r \mathbf{\Sigma}_r \mathbf{V}_r^\top, \quad \mathbf{U}_r \in \mathbb{R}^{d_{\text{out}} \times r}, \; \mathbf{\Sigma}_r \in \mathbb{R}^{r \times r}, \; \mathbf{V}_r \in \mathbb{R}^{d_{\text{in}} \times r}
\end{equation}
سپس تانسور باقیمانده $\mathbf{R}$ از اختلاف وزن متراکم و تقریب رتبه پایین محاسبه می‌گردد:
\begin{equation}
\mathbf{R} = \mathbf{W} - \mathbf{L}
\end{equation}
بر روی باقیمانده $\mathbf{R}$، عملگر هرس برداری ۲:۴ اعمال می‌گردد. در هر بلوک متوالی از ۴ درایه سطری $\mathbf{b} = [r_1, r_2, r_3, r_4]$، دو درایه با کمترین قدرمطلق به صفر تبدیل می‌شوند:
\begin{equation}
\mathbf{R}_{\text{sparse}} = \mathcal{M}_{2:4}(\mathbf{R}), \quad \|\mathbf{b} \odot (1 - \mathbf{m})\|_0 = 2, \quad \mathbf{m} = \arg\max_{\mathbf{m}' \in \{0, 1\}^4, \|\mathbf{m}'\|_1=2} \sum_{k=1}^4 m'_k |r_k|
\end{equation}
ماتریس نهایی بازسازی‌شده از جمع مؤلفه رتبه پایین و باقیمانده پراکنده به دست می‌آید:
\begin{equation}
\mathbf{W}_{\text{recon}} = \mathbf{L} + \mathbf{R}_{\text{sparse}}
\end{equation}
خطای نسبی بازسازی فروبنیوس به عنوان سنجه اعتبارسنجی لایه ارزیابی می‌شود:
\begin{equation}
\mathcal{E}_{\text{rel}} = \frac{\|\mathbf{W} - \mathbf{W}_{\text{recon}}\|_F}{\|\mathbf{W}\|_F + \epsilon}
\end{equation}


\paragraph{تحلیل طیفی و نامساوی‌های پایداری خطا}
برای بررسی میزان دقت تقریب ماتریس وزن یا نگاشت توجه در ماژول \lr{enlighten\_compressor}، تفاضل ماتریس فاکتورگیری‌شده از نگاشت متراکم اصلی بر اساس نامساوی هافدینگ-فون‌نیومان ارزیابی می‌گردد. فرض کنید $\mathbf{W} \in \mathbb{R}^{d_{\text{in}} \times d_{\text{out}}}$ ماتریس وزن هدف باشد که توسط تقریب کم‌مرتبه یا پراکنده با رتبه موثر $r$ جایگزین شده است. کران بالای خطای بازسازی بر مبنای نرم فروبنیوس به صورت زیر تصریح می‌شود:
\begin{equation}
\label{eq:spectral_bound_enlighten_compressor}
\| \mathbf{W} - \mathbf{W}_{r} \|_F = \sqrt{\sum_{i=r+1}^{\min(d_{\text{in}}, d_{\text{out}})} \sigma_i^2(\mathbf{W})} \le \epsilon \| \mathbf{W} \|_F
\end{equation}
که در آن $\sigma_i(\mathbf{W})$ بیانگر $i$-امین مقدار منفرد مرتب‌شده ماتریس به صورت نزولی است. نسبت حفظ انرژی طیفی $\eta_r = \frac{\sum_{i=1}^r \sigma_i^2}{\sum_{i=1}^{d} \sigma_i^2}$ تضمین می‌کند که با انتخاب آستانه $\tau \ge 0.90$، بیش از نود درصد از پراکندگی واریانس داده‌ها حفظ شده و افت سیگنال در انتشار گرادیان رو به عقب به حداقل ممکن کاهش می‌یابد.

\paragraph{تحلیل مجانبی و پیچیدگی اشتراک حافظه سخت‌افزاری}
از منظر پیچیدگی محاسباتی و انتقال داده‌ها در سلسومراتب حافظه (حافظه سراسری \lr{HBM} به حافظه محلی \lr{SRAM})، عملیات ضرب ماتریسی در این مکانیزم نیازمند تخصیص بافر موقت با ابعاد فاکتورگیری‌شده است. برای پردازش دنباله‌ای به طول $N$ با بعد مخفی $D$، بار محاسباتی هسته از فرم درجه دو به رابطه زیر تقلیل می‌یابد:
\begin{equation}
\label{eq:complexity_enlighten_compressor}
\mathcal{C}_{\text{enlighten\_compressor}} = \mathcal{O}\left( B \cdot N \cdot D \cdot r + B \cdot H \cdot N \cdot d_k \right) \ll \mathcal{O}\left( B \cdot H \cdot N^2 \cdot d_k \right)
\end{equation}
این بهبود مرتبه پیچیدگی، با کاهش چشمگیر حجم تبادل تانسورها در گذرگاه‌های ارتباطی کارت گرافیک، مانع از اشباع پهنای باند حافظه (پدیده \lr{Memory-Bound}) در فرآیند استنتاج طولانی‌مدت می‌گردد.
